\section{Covariance decomposition}

\subsection{Decomposition identity}

Write the detrended gap as
\[
e_n = \mu(r_n) + u_n,\qquad \mathbb{E}(u_n \mid r_n) = 0,
\]
where $u_n$ is the within-phase residual. The lag-one covariance decomposes as
\[
\operatorname{Cov}(e_n,\,e_{n+1}) = A + X + Y + D,
\]
with
\[
A = \mathbb{E}\!\left[\mu(r_n)\,\mu(r_{n+1})\right],
\]
\[
X = \mathbb{E}\!\left[\mu(r_n)\,u_{n+1}\right],
\]
\[
Y = \mathbb{E}\!\left[u_n\,\mu(r_{n+1})\right],
\]
\[
D = \mathbb{E}\!\left[u_n\,u_{n+1}\right].
\]
All four terms are computed from the same data; the identity is exact by construction, not an approximation.

\subsection{Values}

The observed components, at $10^{10}$ and 6 decimal places:

We also report the non-$D$ share of the lag-one covariance,
\begin{equation}
  S(M) = \frac{A(M)+X(M)+Y(M)}{\operatorname{Cov}(e_n,e_{n+1})}
       = 1 - \frac{D(M)}{\operatorname{Cov}(e_n,e_{n+1})}.
\end{equation}
Because $\operatorname{Cov}$ and $D$ are both negative, $S(M)$ is
positive and increases with $M$. It is an algebraic decomposition
share, not a causal fraction of dependence removed by conditioning.

\begin{table}[ht]
\centering
\caption{Covariance decomposition components. The sum
$A+X+Y+D$ equals $\operatorname{Cov}$ exactly for every $M$
and is not repeated as a column. The last column reports
$S(M) = (A+X+Y)/\operatorname{Cov}$, the non-$D$ share.}
\label{tab:abcd}
\small
\begin{tabular}{rrrrrrrr}
\toprule
$M$  &  $A$  &  $X$  &  $Y$  &  $X+Y$  &  $D$  &  Cov  &  $S(M)$  \\
\midrule
6 & $-0.014798$ & $-0.011110$ & $-0.142602$ & $-0.153712$ & $-10.832660$ & $-11.001171$ & $1.5\%$ \\
30 & $-0.086958$ & $-0.118493$ & $-1.170132$ & $-1.288625$ & $-9.625587$ & $-11.001171$ & $12.5\%$ \\
210 & $+0.138938$ & $-0.055117$ & $-3.084120$ & $-3.139237$ & $-8.000871$ & $-11.001171$ & $27.3\%$ \\
2310 & $+0.670972$ & $-0.035007$ & $-4.876237$ & $-4.911244$ & $-6.760899$ & $-11.001171$ & $38.5\%$ \\
30030 & $+1.369692$ & $-0.029422$ & $-6.575378$ & $-6.604800$ & $-5.766063$ & $-11.001171$ & $47.6\%$ \\
510510 & $+2.127980$ & $-0.028428$ & $-8.085351$ & $-8.113779$ & $-5.015372$ & $-11.001171$ & $54.4\%$ \\
9699690 & $+2.898121$ & $-0.031791$ & $-9.515775$ & $-9.547566$ & $-4.351726$ & $-11.001171$ & $60.4\%$ \\
\bottomrule
\end{tabular}
\end{table}

\subsection{Qualitative features}

Three features stand out.

\textbf{(i) Sign change of $A$.} The phase component $A$ is small and negative for $M=6$ and $M=30$, then changes sign between $M=30$ and $M=210$, and grows monotonically for $M\ge 210$, reaching $2.898$ at $M=9699690$.

\textbf{(ii) Dominance of $D$.} The residual term $D$ dominates the decomposition at small $M$ ($D=-10.833$ out of $\operatorname{Cov}=-11.001$ at $M=6$), and decreases in magnitude as $M$ grows, reaching $-4.352$ at $M=9699690$. The residual remains negative at every tested modulus.

\textbf{(iii) Co-dominance of $Y$ and $D$ at large $M$.} At small $M$ ($M=6$, $30$) the residual term $D$ dominates the decomposition. At larger $M$ the cross term $Y$ grows in magnitude and becomes comparable to $D$: at $M=9699690$, $Y = -9.516$ while $D = -4.352$. The two co-dominate at the largest modulus tested. The term $X$ remains small at every $M$. The phase component $A$ changes sign between $M=30$ and $M=210$ and grows monotonically for $M\ge 210$.

\subsection{Residual correlation}

Define the residual lag-one correlation
\[
\rho_{\text{resid}}(M) = \frac{D(M)}{\text{var\_within}(M)}.
\]
The observed values are reported in Table~\ref{tab:rho-resid}:

\begin{table}[ht]
\centering
\caption{Residual lag-one correlation.}
\label{tab:rho-resid}
\begin{tabular}{rr}
\toprule
$M$ & $\rho_{\text{resid}}$ \\
\midrule
6 & $-0.024703$ \\
30 & $-0.022007$ \\
210 & $-0.018402$ \\
2310 & $-0.015653$ \\
30030 & $-0.013445$ \\
510510 & $-0.011775$ \\
9699690 & $-0.010306$ \\
\bottomrule
\end{tabular}
\end{table}

The residual correlation remains negative at all seven tested scales. Its magnitude declines from $0.0247$ at $M=6$ to $0.0103$ at $M=9699690$ --- a reduction of roughly $58\%$. Phase conditioning therefore reduces the residual term $D$ from $-10.83$ at $M=6$ to $-4.35$ at $M=9699690$, while the total covariance $A+X+Y+D$ remains fixed at $-11.001171$. The redistribution among the four components does not eliminate the residual dependence within the window.

The CSV also contains a small-sample-corrected variant $\rho_{\text{resid}}^{\text{corr}}(M)$, which differs from the raw value only at the largest modulus: $-0.010326$ versus $-0.010306$ at $M=9699690$. The correction is negligible at all tested scales and the paper reports the raw values.

\subsection{Block bootstrap}

Confidence intervals for $\rho_{\text{resid}}$ are obtained by a moving-block bootstrap with 40 contiguous blocks and 1000 replicates, with an independent seed for each $M$. The $95\%$ intervals are reported in Table~\ref{tab:bootstrap}:

\begin{table}[ht]
\centering
\caption{Block-bootstrap 95\% confidence intervals for $\rho_{\text{resid}}$.}
\label{tab:bootstrap}
\begin{tabular}{rrrr}
\toprule
$M$ & $\rho_{\text{resid}}$ & 95\% CI & width \\
\midrule
6 & $-0.024703$ & $[-0.024868, -0.024544]$ & 0.000324 \\
30 & $-0.022007$ & $[-0.022138, -0.021864]$ & 0.000274 \\
210 & $-0.018402$ & $[-0.018511, -0.018283]$ & 0.000228 \\
2310 & $-0.015653$ & $[-0.015750, -0.015551]$ & 0.000199 \\
30030 & $-0.013445$ & $[-0.013535, -0.013352]$ & 0.000183 \\
510510 & $-0.011775$ & $[-0.011856, -0.011691]$ & 0.000165 \\
9699690 & $-0.010306$ & $[-0.010379, -0.010226]$ & 0.000153 \\
\bottomrule
\end{tabular}
\end{table}

All seven marginal 95\% bootstrap intervals exclude zero. The point estimates also decrease monotonically across the tested modulus ladder. We do not claim a familywise test of monotonicity.

\subsection{What this section does not claim}

$A$, $X$, $Y$, $D$ are conditional components computed from the same data. The decomposition does not establish a causal role for the phase component. It establishes that, under the present conditioning, the residual term $D$ remains negative at all seven tested moduli, and that the total covariance is redistributed across $A$, $X$, $Y$, $D$ as $M$ grows: $A$ becomes increasingly positive, $Y$ becomes increasingly negative, and $D$ decreases in magnitude. The sum $A+X+Y+D$ is unchanged.

No claim is made that $\rho_{\text{resid}}(M)\to 0$ as $M\to\infty$. The seven moduli are a finite ladder; extrapolation beyond $M=9699690$ is not supported by the present data.

A note on the mechanical origin of $Y$. The next phase class satisfies $r_{n+1} = r_n + g_n \bmod M$, so the gap $g_n$ determines $r_{n+1}$ exactly. Part of the transfer of covariance from $D$ to $Y$ as $M$ grows is therefore a re-distribution within the accounting identity, not a physical removal of dependence. The total $A+X+Y+D$ is fixed at $-11.001171$ for every $M$; what changes is its distribution across the four components. This identity is also the reason that $Y$ is conceptually related to the consecutive-residue transition viewpoint discussed in Section~\ref{sec:theory}.
